\section{Bewertung} \label{sec:bewertung_optimierung}
Die Bewertung der vorgenommenen Optimierungen des Gipfelfindungsalgorithmus erfolgt anhand von Laufzeitmessungen und der Analyse der Ergebnisse des Algorithmus. Es wurden erneut die in \ref{sec:testen_des_algorithmus} beschriebenen Testfälle durchgeführt, um die Auswirkungen der Optimierungen auf die Genauigkeit der Ergebnisse zu überprüfen. Dabei lag der Fokus insbesondere auf der Prominenz- und Dominanzabweichung und der Anzahl der gefundenen realen Gipfel, sowie den fehlklassifizierten Gipfeln. Die Optimierungen haben in allen Testfällen zu einer Verbesserung der Genauigkeit des Algorithmus geführt. ein gutes Beispiel dafürr ist der Wettersteintest. Bei diesem werden nun keine nicht existierenden Gipfel mehr gefunden und fast alle realen Gipfel. Als einziger realer Gipfel wird der kleine Solstein nicht korrekt erkannt, da seine höhe in der GeoTIFF Datei neidriger angegeben ist und er dadruch nicht mehr die Prominenzbedingung des Tests erfüllt. Außerdem haben sich die Prominenz-, sowie Dominanzabweichung stark verringert. In dem Wettersteintest ist durch die Optimierungen die Prominenzabweichung und Dominanzabweichung um ca. Faktor 4 geringer als zuvor. Bei den anderen Tests liegt die Verringerung der Abweichungen immer über dem Faktor 2 und im Mittel bei 2,75. Die Laufzeit hat sich auch verbessert. Diese wurden in \ref{sec:laufzeitverbesserung:optimierung} thematisiert. 
